%=====================================================================
% 機械学習 第1回 演習問題 提出用テンプレート
%
% 使い方
%   1. 下の \学籍番号, \氏名, \提出日 を自分の情報に書き換える
%   2. 各問題の「解答」枠の中に解答を書く（枠の外の問題文は変更しない）
%   3. コンパイルして PDF を作成し、PDF を提出する
%
% コンパイル方法（ターミナル）
%   uplatex report_template_01.tex && dvipdfmx report_template_01.dvi
%
% Overleaf を使う場合
%   Menu > Compiler を "LaTeX" にし、プロジェクトに latexmkrc という名前の
%   ファイルを作って次の3行を書く：
%     $latex = 'uplatex';
%     $dvipdf = 'dvipdfmx %O -o %D %S';
%     $pdf_mode = 3;
%
% ファイル名は  ML01_学籍番号.pdf  として提出すること（例：ML01_2312345.pdf）
%=====================================================================
\documentclass[a4paper,dvipdfmx]{jsarticle}
\usepackage{amsmath}
\usepackage{amssymb}
\usepackage{bm}
\usepackage{ascmac}
\usepackage[dvipdfmx]{hyperref}
\usepackage{pxjahyper}
\hypersetup{pdfborder={0 0 0}}

%---- 自分の情報を書く -------------------------------------------------
\newcommand{\学籍番号}{0000000}
\newcommand{\氏名}{氏名を書く}
\newcommand{\提出日}{2026年10月1日}
%----------------------------------------------------------------------

% 解答枠（枠の中に解答を書く）
\newenvironment{answer}{\begin{itembox}[l]{解答}}{\end{itembox}}

\begin{document}

\begin{center}
  {\LARGE 機械学習　第1回 演習問題}\\[2mm]
  {\large 提出課題}\\[6mm]
  \begin{tabular}{ll}
    学籍番号： & \学籍番号 \\
    氏名：     & \氏名 \\
    提出日：   & \提出日 \\
  \end{tabular}
\end{center}

\vspace{4mm}
\noindent
\textbf{注意}
\begin{itemize}
  \item 解答は「解答」の枠内に書くこと。枠外の問題文は変更しないこと。
  \item 計算問題は結果だけでなく途中の式も書くこと。
  \item 判定問題（正しい／誤り）は理由を必ず書くこと。
  \item 数式の書き方は末尾の「LaTeX の書き方」を参照すること。
\end{itemize}

\section*{問題}

\begin{enumerate}
  \item \textbf{記法の確認}\quad
    次の各主張が正しいか否かを判定し、その理由を述べよ。
    \begin{enumerate}
      \item $\bm{x} \in \mathbb{R}^3$ かつ $\mathcal{X} \subseteq \mathbb{R}^3$ であれば、
        常に $\bm{x} \in \mathcal{X}$ が成り立つ。
      \item $h : \mathbb{R}^2 \to \{0, 1\}$ は、2次元入力に対する2値分類の仮説を表す。
      \item $\mathcal{Y} = \{0, 1\}$ のとき、直積 $\mathcal{Y} \times \mathcal{Y}$ の要素の個数は $4$ である。
      \item $h(\cdot\,; \bm{w})$ と $h(\bm{x}; \bm{w})$ は、どちらも関数そのものを指す記法である。
    \end{enumerate}
    \begin{answer}
      \begin{enumerate}
        \item （正しい／誤り）理由：
        \item （正しい／誤り）理由：
        \item （正しい／誤り）理由：
        \item （正しい／誤り）理由：
      \end{enumerate}
    \end{answer}

  \item \textbf{パラメータ次元}\quad
    次の各仮説空間のパラメータ次元 $p$ を求めよ。
    \begin{enumerate}
      \item 入力 $\bm{x} \in \mathbb{R}^5$ に対する線形モデル $h(\bm{x}; \bm{w}) = w_0 + \bm{w}_{1:5}^\top \bm{x}$
      \item 入力 $x \in \mathbb{R}$ に対する3次多項式モデル
      \item 入力 $\bm{x} = (x_1, x_2)^\top \in \mathbb{R}^2$ に対して、
        2次以下のすべての単項式 $1, x_1, x_2, x_1^2, x_1 x_2, x_2^2$ の線形結合で表されるモデル
    \end{enumerate}
    \begin{answer}
      \begin{enumerate}
        \item $p = $
        \item $p = $
        \item $p = $
      \end{enumerate}
    \end{answer}

  \item \textbf{仮説空間の包含関係}\quad
    $\mathcal{H}_M$ を $M$ 次多項式仮説空間とする。
    $\mathcal{H}_1 \subset \mathcal{H}_2$（$\mathcal{H}_1$ は $\mathcal{H}_2$ の真部分集合）であることを、
    次の2点を示すことで証明せよ。
    \begin{enumerate}
      \item 任意の $h \in \mathcal{H}_1$ は $\mathcal{H}_2$ の要素でもある。
      \item $\mathcal{H}_2$ の要素であって $\mathcal{H}_1$ の要素ではない関数が存在する。
    \end{enumerate}
    \begin{answer}
      \begin{enumerate}
        \item
        \item
      \end{enumerate}
    \end{answer}

  \item \textbf{経験損失の計算}\quad
    訓練データ $D = \{(1, 3),\; (2, 5),\; (3, 6)\}$ が与えられている
    （各組は $(x_i, y_i)$ を表す）。
    2つの仮説
    \begin{equation*}
      h_a(x) = 2x + 1, \qquad h_b(x) = 1.5x + 1.5
    \end{equation*}
    について、二乗損失に基づく経験損失 $L_D(h_a)$, $L_D(h_b)$ をそれぞれ計算し、
    どちらの仮説が訓練データに対して良いかを判定せよ。
    \begin{answer}
      \begin{align*}
        L_D(h_a) &= \\
        L_D(h_b) &=
      \end{align*}
      判定：
    \end{answer}

  \item \textbf{0-1損失}\quad
    5個の訓練データに対する真のラベルと仮説 $h$ の予測が次のように与えられている。
    \begin{center}
      \begin{tabular}{c|ccccc}
        $i$ & 1 & 2 & 3 & 4 & 5 \\
        \hline
        $y_i$ & 1 & 0 & 1 & 1 & 0 \\
        $h(\bm{x}_i)$ & 1 & 0 & 0 & 1 & 1 \\
      \end{tabular}
    \end{center}
    0-1損失に基づく経験損失 $L_D(h)$ を求めよ。
    また、この値が一般に何と呼ばれる量に一致するかを述べよ。
    \begin{answer}
      \begin{equation*}
        L_D(h) =
      \end{equation*}
      この値は（　　　　）に一致する。
    \end{answer}

  \item \textbf{$\min$ と $\arg\min$}\quad
    問4の設定で、仮説空間を $\mathcal{H} = \{h_a, h_b\}$ の2要素からなる集合とする。
    このとき
    \begin{equation*}
      \min_{h \in \mathcal{H}} L_D(h), \qquad h^* = \arg\min_{h \in \mathcal{H}} L_D(h)
    \end{equation*}
    をそれぞれ求めよ。両者の違いを一言で説明せよ。
    \begin{answer}
      \begin{align*}
        \min_{h \in \mathcal{H}} L_D(h) &= \\
        h^* = \arg\min_{h \in \mathcal{H}} L_D(h) &=
      \end{align*}
      両者の違い：
    \end{answer}

  \item \textbf{経験損失と汎化誤差}\quad
    平面上の3点 $(x_1, y_1), (x_2, y_2), (x_3, y_3)$（$x_i$ は互いに異なる）が与えられたとき、
    これらをすべて通る2次多項式 $h \in \mathcal{H}_2$ がただ一つ存在する。
    \begin{enumerate}
      \item この $h$ の二乗損失に基づく経験損失 $L_D(h)$ を求めよ。
      \item この $h$ の汎化誤差 $L_{\mathrm{gen}}(h)$ も $0$ であると言えるか。
        理由とともに述べよ。
      \item 一般に $N$ 個のデータ点（$x_i$ は互いに異なる）が与えられたとき、
        経験損失を $0$ にできる多項式の次数 $M$ の最小値はいくらか。
        この事実と過学習の関係を説明せよ。
    \end{enumerate}
    \begin{answer}
      \begin{enumerate}
        \item
        \item
        \item
      \end{enumerate}
    \end{answer}
\end{enumerate}

\clearpage
\section*{LaTeX の書き方（参考）}

このテンプレートで使う数式の書き方をまとめる。
数式は \verb|$ ... $| で囲む（行の中に書く場合）か、
\verb|\begin{equation*} ... \end{equation*}| で囲む（1行独立させる場合）。

\begin{center}
\small
\begin{tabular}{l|l}
\hline
表示 & 書き方 \\
\hline\hline
$\bm{x}$（ベクトル） & \verb|\bm{x}| \\
$\mathcal{H}$, $\mathcal{X}$（集合） & \verb|\mathcal{H}|, \verb|\mathcal{X}| \\
$\mathbb{R}$, $\mathbb{R}^3$ & \verb|\mathbb{R}|, \verb|\mathbb{R}^3| \\
$\bm{x} \in \mathcal{X}$ & \verb|\bm{x} \in \mathcal{X}| \\
$\bm{x} \notin \mathcal{X}$ & \verb|\bm{x} \notin \mathcal{X}| \\
$\mathcal{X} \subseteq \mathbb{R}^3$, $\mathcal{H}_1 \subset \mathcal{H}_2$ & \verb|\mathcal{X} \subseteq \mathbb{R}^3|, \verb|\mathcal{H}_1 \subset \mathcal{H}_2| \\
$h : \mathcal{X} \to \mathcal{Y}$ & \verb|h : \mathcal{X} \to \mathcal{Y}| \\
$\{0, 1\}$, $\mathcal{Y} \times \mathcal{Y}$ & \verb|\{0, 1\}|, \verb|\mathcal{Y} \times \mathcal{Y}| \\
$\{\, h \mid h \in \mathcal{H} \,\}$ & \verb|\{\, h \mid h \in \mathcal{H} \,\}| \\
$\bm{x}^\top$（転置） & \verb|\bm{x}^\top| \\
$w_0$, $x_1^2$（添字・べき） & \verb|w_0|, \verb|x_1^2| \\
$h(\bm{x}; \bm{w})$, $h(\cdot\,; \bm{w})$ & \verb|h(\bm{x}; \bm{w})|, \verb|h(\cdot\,; \bm{w})| \\
$\dfrac{1}{N}$（分数） & \verb|\frac{1}{N}| \\
$\displaystyle\sum_{i=1}^{N} (y_i - h(x_i))^2$ & \verb|\sum_{i=1}^{N} (y_i - h(x_i))^2| \\
$L_D(h)$, $L_{\mathrm{gen}}(h)$ & \verb|L_D(h)|, \verb|L_{\mathrm{gen}}(h)| \\
$\displaystyle\min_{h \in \mathcal{H}} L_D(h)$ & \verb|\min_{h \in \mathcal{H}} L_D(h)| \\
$\displaystyle h^* = \arg\min_{h \in \mathcal{H}} L_D(h)$ & \verb|h^* = \arg\min_{h \in \mathcal{H}} L_D(h)| \\
$a \neq b$, $a \leq b$, $a \geq b$ & \verb|a \neq b|, \verb|a \leq b|, \verb|a \geq b| \\
$\cdots$, $\ldots$ & \verb|\cdots|, \verb|\ldots| \\
\hline
\end{tabular}
\end{center}

\noindent
\textbf{計算過程を複数行で書く例}（\verb|align*| 環境。\verb|&| の位置で揃い、\verb|\\| で改行する）：
\begin{verbatim}
\begin{align*}
  L_D(h) &= \frac{1}{3} \sum_{i=1}^{3} (y_i - h(x_i))^2 \\
         &= \frac{(3-3)^2 + (5-5)^2 + (6-7)^2}{3} \\
         &= \frac{1}{3}
\end{align*}
\end{verbatim}
表示結果：
\begin{align*}
  L_D(h) &= \frac{1}{3} \sum_{i=1}^{3} (y_i - h(x_i))^2 \\
         &= \frac{(3-3)^2 + (5-5)^2 + (6-7)^2}{3} \\
         &= \frac{1}{3}
\end{align*}

\noindent
\textbf{場合分けの例}：
\begin{verbatim}
\begin{equation*}
  L(y, h(\bm{x})) =
  \begin{cases}
    0 & (y = h(\bm{x})) \\
    1 & (y \neq h(\bm{x}))
  \end{cases}
\end{equation*}
\end{verbatim}
表示結果：
\begin{equation*}
  L(y, h(\bm{x})) =
  \begin{cases}
    0 & (y = h(\bm{x})) \\
    1 & (y \neq h(\bm{x}))
  \end{cases}
\end{equation*}

\end{document}
